% Project Proposal
% Total points: 100 
\documentclass[11pt]{article}
\usepackage[margin=1in]{geometry}
\usepackage{hyperref}
\title{Proposal on Success of Civil War Recovery}
\author{Hazel Lutz}
\date{January 21, 2026}

\begin{document}
\maketitle

\section{Objective (5 points)}

After over a decade, the Syrian civil war came to a (tentative) close in December of 2024, when rebel leader Ahmed al-Sharaa and his rebel group Hay’at Tahrir al-Sham stormed Syria’s capital, leading to former dictator Bashar al-Assad’s flight from the city. In the year since, the new Syrian government has engaged in numerous measures to re-establish sovereignty and legitimacy in a country still fractured by numerous rebel groups and regime elements, each with a degree of external support. Continued conflict between the current government and other groups with claims to sovereignty, such as the Kurdish-led Syrian Democratic Forces after a breakdown in transition negotiations and the Israeli-backed Druze in the southwest, threaten to destabilize reconstruction efforts and cause a resurgence of war. Understanding the necessary conditions for reconstruction in Syria, especially as it relates to historical cases of post-conflict reconstruction, is therefore necessary.

In the current low-data environment Syria presents, understanding its needs with regards to post-conflict reconstruction will require learning from past cases. To do so, I will use Bayesian methods to model how the success of recovery (and potentially duration?) are impacted by variables related to the conflict itself (extent of external support for armed factions, number of armed factions, infrastructural damages, cause of war, rebel or government victory) and variables related to pre-conflict conditions (legitimacy of previous state, economic disparity, etc.). This will conclude with a brief comparative analysis between in-region cases that may prove relevant in contextualizing the quantitative analysis, such as the Lebanese Civil War from 1975-1990, the North Yemen civil war from 1962-1970, or the ongoing Yemeni civil war.


\section{Outcome (5 point)}

The outcomes of this will be the creation of an updating (using APIs) database of previous civil wars (time range TBD, but worth restricting due to the changing nature of civil wars in recent decades as well as public opinion data availability) and the creation of different models for different types of civil wars, whether protracted or involving external actors tipping the scales in their favor, to account for differences between cases like Syria’s or unlike Syria’s (demonstrating whether cases like these are truly unique).

\section{Impact (5 points)}

There are two goals for this project: first, to provide empirical backing for a theoretical distinction between civil wars based on external involvement due to differential outcomes for reconstruction, and second, to provide direction on the necessary conditions for reconstruction in Syria’s case compared to others. The former is of most interest to political science researchers seeking to understand the applicability of past literature to modern cases – in which modern technology has both created powerful security states, such as Syria under Bashar al-Assad, as well as altered the nature of conflict itself. The latter impact is most important for those interested in reconstruction policy, both in aid organizations and policy research groups – improving the on-the-ground circumstances in Syria and similar war-torn states remains the guiding star of this study.

\section{Literature Review (25 points)}
% Do a thorough literature review to explore what has already been studied in the direction of your research question. Cite at least five most relevant publications below, summarizing their contributions. 
\begin{enumerate}
    \item \textbf{Sharif, "Predicting the End of the Syrian Conflict: From Theory to the Reality of a Civil War" (2021):} Sharif's article will be foundational for understanding distinct features to look for when considering the particular kind of civil war -- perhaps a less traditional kind -- that I am interested in studying. In this article, she considers incorrect predictions about the length and intensity of the Syrian civil war, whose inaccuracies result from "the problematic cartography of ongoing conflicts, the splintering of the opposition..., as well as disagreement among foreign interveners" (p. 327). This, she argues, makes it elude standard categorizations. Stripped down, my fundamental research question is similar: to what extent are categorizations of civil war \emph{recovery} insufficient in dealing with cases like the Syrian civil war? What special considerations must be made?
    \item \textbf{Doyble and Sambanis, "International Peacebuilding: A Theoretical and Quantitative Analysis" (2000)}: Using a dataset they have constructed of civil wars between 1944 and 1997 (p. 783), Doyle and Sambanis test a number of hypotheses related to the success of peacebuilding, providing a number of variables to include in regression (length, cause of war, economi situation pre-war, etc.) that were tested in their analysis (pp. 783-786). Additionally, they give definitional criteria that should be considered, in length, number of deaths, boundaries, etc (p. 783). A greater interrogation of the methodology of this article will be necessary to derive some lessons for my own analysis. The conclusions of this provide great support to theories about the efficacy of international intervention and what kinds of peacebuilding ("higher order" vs. "lower order") can be expected (p. 795).
    \item \textbf{Sisk, "Power-Sharing in Civil War: Puzzles of Peacemaking and Peacebuilidng" (2013)}: In the wake of the Syrian civil war, this article considers the efficacy of power-sharing agreements and different models of externally-influenced peace. Written over a decade before the conclusion of the war, the actual outcome certainly could not have been anticipated, but this provides some theoretical underpinnings for traditional and alternative methods of peacebuilding. Syria's current trajectory seems to be one that rejects this sort of power-sharing model, at least in the most rigid, externally-imposed sense between equal powers, such as the "confessional" system after the Lebanese civil war (p. 8); whether this is a beneficial trajectory is worth considering.
    \item \textbf{Dixon, "What Causes Civil Wars? Integrating Quantitative Research Findings" (2009)}: In addition to being beneficial to consider methodologically in providing quantitative studies to look at, this includes many considerations for pre-war/war-initiating circumstances that may continue to be influential in recovery, such as the kind of grievance that caused the civil war, characteristics of the population, economic issues, etc. When building a model, the features tested here will be very useful; however, it provides little direction for the theoretical grounding of what I will be studying.
    \item \textbf{Adedokun, "Peace Building After Civil War: A Critical Survey of the Literature and Avenues for Future Research" (2017)}: This article will be useful for some of the theory that will underlie my quantitative analysis -- fundamentally, while many quantitative studies, especially using regression, consider a varied number of different determinants for some outcome (e.g., successful postwar peacebuilding), this covers that literature on theoretical grounds on a meta-analysis level and finds directions for future research. While the meat of the paper, which performs this meta-analysis, will give me a broad overview of both peace studies and conflict studies literature and the general trends in it, the conclusion especially provides a few prescriptions: namely, attention to "large N studies on why peacebuilding sometimes succeeds", lacking understanding on "what constitutes local ownership" (bottom-up, integrated peace-building approaches rather than top-down, imposed ones), "developing appropriate criteria for... evaluating what comprises 'success' or 'failure'", and "the integration of qualitative and quantitative approaches" (p. 42). All four of these considerations will be essential to my analysis: while some provide questions for theoretical clarification necessary for adequate quantitative discussions, others are incorporated into my methodology and goals. I look forward to contributing to some of the gaps explored here.
\end{enumerate}

\section{Novelty (10 points)}
Though detailed to some degree in my literature review itself, many of the areas surveyed provide answers to some tangential questions, focusing on slightly different cases or attributes than I have, and others provide useful theoretical elaborations that, in some cases, explicitly ask for broader quantitative research to be performed. My research has been shaped by these gaps I have found in the literature, and will likely continue to be as my research goes onward (though some attention to retaining focus and scope will be necessary). And while some studies of post-war reconstruction have been done, mine will not be studying what is correlated with success in post-war reconstruction, but rather what necessities for reconstruction exist for different kinds of wars.   

\section{Data Source(s) (10 point)}
\begin{itemize}
    \item \href{https://ucdp.uu.se/downloads/}{UCDP GED}: for conflict data, generally very reliable, includes geographic location and type of conflict
    \item \href{https://data.worldbank.org/}{World Bank}: useful for economy, infrastructure, development, though some data may be sparse
    \item \href{https://www.v-dem.net/data/}{V-Dem}: useful for surrounding variables regarding social conditions, and data is independently determined and coded, so generally very available
    \item Public opinion (\href{https://www.vanderbilt.edu/lapop/raw-data.php}{AmericasBarometer}, \href{https://www.arabbarometer.org/survey-data/data-downloads/}{Arab Barometer}, \href{https://www.afrobarometer.org/data/merged-data/}{Afrobarometer}, etc.): available countries and questions asked may differ by year, requiring a level of compromising on needs and data availability
\end{itemize}

\section{Approach (20 points)}

The first step to this will be creating a dataset of civil wars and relevant independent variables – this will involve, first, using a database of conflicts and categorizing them based on the nature of the conflict (based on external involvement and duration), then sourcing and aggregating data for other relevant, aforementioned variables about the nature of the conflict and the state itself (from the World Bank, V-Dem, etc?). This is likely to be somewhat lengthy due to a degree of manual work, research on variables and suitable proxy measures, and the nature of the available data will heavily alter the analysis performed.

The next main step will be the start of  Bayesian modeling – this requires not only the selection of suitable priors (with testing for the sensitivity of models based on the prior selected), but also determining the nature of the Bayesian model based on the data collected. In anticipation of a small/medium N dataset, Bayesian methods are what I believe would be most useful, and while Bayesian multiple regression may well be sufficient, more complex methods may prove necessary. In preparation of this, I am surveying Bayesian social scientific methods at the same time as creating my dataset – a more in-depth summary of potential methods will be present in the final version of this proposal.

Following this, I will perform a comparative political analysis of illustrative cases with a focus on providing the qualitative backdrop of the variables used in my Bayesian model to understand how they interact with circumstances on the ground that are harder to quantify or are not as yet operationalized. Cases will be selected based on the preceding analysis, but priority will be given to cases that are within the Middle East if available, such as the Lebanese civil war from 1975-1990. This will be used to give further direction to the application of the data to Syria’s current reconstruction efforts.

\section{Timeline (10 points)}
\begin{itemize}
    \item (4 Weeks): Creating (includes coding cases)/automating database, simultaneous research on methods and important variables
    \item (1 Week): EDA – variable selection, potential multicollinearity, revising grouping, etc.
    \item (4 Weeks): Bayesian Modeling 
    \item (0.5 Weeks): Compiling results
    \item (1.5 Weeks): Contextualizing/qualifying quantitative results with comparative political analysis
    \item (1 Week): Connecting data to Syria’s case
    \item (1 Week): Writing up/revising report
    \item (1 Week): Hosting database
    \item (1 Week): Poster and Final Presentation
\end{itemize}
\section{Possible Challenges (10 points)}

Most issues will come from data availability and methodology, but this is not my first time with such difficulties – frequently, in political science research, datasets will be sparse or not include consistent information due to difficulty in gathering data or limited funding, especially with public opinion data. This requires flexibility and the use of proxy measures as necessary, and Bayesian methods will specifically be useful where data is limited. I do not anticipate many issues elsewhere.


\bibliographystyle{plain}
\bibliography{references}

\end{document}